\section{The Propers: Detailed Commentary}

\subsection{A threatened people before the creator (Introit)}
Mordecai's prayer arises within a concrete danger. Haman has obtained a
decree against the Jews; Mordecai and Esther must act in a royal court
where access to the king is itself perilous. The prayer in Esther~13:8--18
belongs to Mordecai; Esther's own prayer follows in chapter~14. Its claim
that nothing can resist God is therefore a plea from endangered people,
not the voice of an earthly ruler announcing unlimited power. Heaven,
earth, and everything beneath heaven belong to the God whom they invoke.
Haman's power is formidable, but it is not ultimate.\footnote{Esth.~3--4;
13:8--18;14, Douay--Rheims--Challoner. Esther~13:9--11 belongs to the Greek
additions, corresponding to Addition C:2--4; chapter numbers vary across
editions.}

The chant concentrates that confession. The Bible's conditional phrase
about God's determination to save Israel is absent, while the address
\textit{Domine} and the concluding declaration of universal lordship give
the antiphon a direct liturgical force. Its psalm verse then changes the
angle: those who invoke the Lord must walk in his law. The undefiled way
is not a claim that the worshippers need no mercy. Augustine observes that
the same psalm soon wishes for its ways to be directed by God. A wish for
help reveals dependence; an order to walk does not make the walker
self-sufficient.\footnote{Augustine, \work{Expositions of the Psalms},
Ps.~118, Aleph \S\S1--8, especially 3,5--7; NPNF I.8, under Ps.~119.}

\subsection{Freedom for good works (Sunday Collect)}
\textit{Familiam tuam} makes the people God's household before it asks
what they should do. The prayer moves from continuing mercy, through
protection and freedom from adversity, to devotion expressed in good
acts. It does not ask for freedom as an empty space in which no claim can
be made on the family. Freedom has a purpose: serving God's name. The
protecting subject is God throughout; the family receives and acts.

Schuster hears this communal dimension strongly. The Christian does not
ask for private security detached from the Church's life. Protection
overcomes obstacles to righteous action and opens into devoted service.
His account helps connect the plural household with the Epistle's common
armor. Yet the Collect's request remains concrete: it really asks release
from adversities. The later appearance of wounded Job will deepen that
petition's horizon without changing it into a request to be
wounded.\footnote{Ildefonso Schuster, \work{The Sacramentary}, trans.
Arthur Levelis-Marke, III (1927), p.~179.}

\subsection{Armor for peace (Epistle)}
The armor concludes Ephesians' sustained teaching about a renewed life.
Anger must not become sin; bitterness and malice yield to forgiveness;
the community imitates the love already shown it in Christ. Immediately
before the appointed passage, earthly masters are reminded that they
and their servants have the same Master in heaven. The strength sought
in verse~10 belongs to that shared Lord. It is no endorsement of one
person's unlimited power over another.\footnote{Eph.~4:26--32;5:1--2;
6:5--10. Chrysostom, \work{Homilies on Ephesians} XXII, on 6:9--13,
NPNF I.13.}

The repeated command is to stand. Truth binds together the person who
would otherwise be pulled apart by falsehood. Justice protects the
breast; peace prepares the feet; faith meets fiery darts; salvation guards
the head; the Spirit supplies the word of God as sword. Each object is
explained by the virtue or gift named with it. The metaphor does not
authorize a corresponding set of material weapons. Its first exclusion,
``not against flesh and blood,'' governs everything that follows.

Chrysostom understands the breastplate as the whole virtuous life and
the Gospel of peace as reconciliation with God. He gives more than one
account of the footwear: readiness to serve the Gospel, and readiness for
departure. He initially treats the shield's faith as a gift of power,
then develops trust in the coming good that withstands doubt and
distress. His detailed exposition is wider than a fixed list of
one-to-one symbols.\footnote{Chrysostom, \work{Ephesians} XXIV, on
6:14--17 and the subsequent moral exposition.}

Two differences deserve precision. Chrysostom takes ``in the heavenly
places'' to name the heavenly goods at stake in the contest. The NPNF
editor instead defends a locative construction. Chrysostom's ``evil
day'' is the present life, whereas the edition's note prefers the day
of violent temptation and also discusses a final outbreak. All can
prompt vigilance, but their explanations of the words are not identical.
The Mass appoints verses~10--17; the prayer for all the saints in
verse~18 remains their immediate scriptural context.\footnote{Chrysostom,
\work{Ephesians} XXII; NPNF I.13, editorial notes 465 and 468.}

\subsection{Refuge before the mountains (Gradual)}
Psalm~89 begins with Moses in its title and moves from God's eternity
to human mortality. The Gradual selects the opening assurance. Mountains
seem older and stronger than a human life, yet they too were made. The
Lord who precedes them is the refuge of successive generations. The
assembly's danger is not made unreal; it is set beside a fidelity that
does not begin or end with this generation.

Augustine follows the tense of ``thou art.'' God is not merely someone
who was strong long ago and will perhaps be strong again. His being is
unchanging. God becomes our refuge when we turn to him without thereby
beginning to exist or undergoing a change in himself. Bellarmine similarly
reads the prayer's memory of past help as reason for present confidence:
faith, hope, and charity make the divine refuge a lived dwelling.
Augustine expressly declines to infer Mosaic composition simply from
the heading.\footnote{Augustine, Ps.~89 \S\S1--4 (NPNF Ps.~90);
Robert Bellarmine, \work{A Commentary on the Book of Psalms}, trans.
John O'Sullivan (1866), Ps.~89:1--3.}

\subsection{The house that leaves bondage (Alleluia)}
The appointed verse is the beginning of Psalm~113, not the whole psalm's
account of sea, Jordan, mountains, and rock. Its sentence carries the
house of Jacob out of Egypt. The historical English ``barbarous people''
marks the foreign people and language from which Israel departed;
Bellarmine explains the language difference. The verse does not establish
an ethnic scale of human worth.\footnote{Bellarmine, Ps.~113:1;
Ps.~113:1--8 (Hebrew Ps.~114).}

Augustine reads the Exodus through Paul's account of the cloud, sea,
spiritual food, and spiritual drink in 1 Corinthians~10. The Church,
gathering Jewish and Gentile believers in Christ, recognizes her own
deliverance there. But Paul's warning belongs to the same argument:
many who passed through the sea subsequently fell. Deliverance creates
a people called to fidelity; it does not make fidelity unnecessary.
This warning will become particularly searching when the pardoned
servant leaves his master's presence.\footnote{Augustine, Ps.~113
\S\S1--6 (NPNF Ps.~114), especially 3--5.}

\subsection{The debt and the repeated plea (Gospel)}
Matthew~18 has already taught concern for little ones, recovery of the
lost, and correction of a brother. Peter then asks how often he must
forgive. The appointed Gospel begins with the king's accounting in
verse~23, after that question. Its dramatic sequence matters: impossible
liability, a plea for time, complete remission, a second plea, refusal,
report, and judgment. The first servant is not condemned for lacking
resources to repay the king. He is condemned after refusing a fellow
servant the very patience he had asked for himself.

Chrysostom draws attention to two unequal exchanges. The king gives more
than delay: he releases the debtor. The debtor gives less than the delay
his companion asks: he casts him into prison. The repetition of the
petition should have recalled his own rescue. Instead, the freedom just
given him becomes the means of throttling another. The final demand for
forgiveness ``from your hearts'' names the interior contradiction behind
the visible violence. The first complete reading below follows this
problem through Chrysostom, Augustine, and Aquinas.\footnote{Chrysostom,
\work{Homilies on Matthew} LXI, \S\S3--5, NPNF I.10; Augustine,
Sermon~83, \S\S1--2,7--8, NPNF I.6, Sermon XXXIII.}

\subsection{The righteous man who is wounded (Offertory)}
The Offertory begins by naming Job upright and God-fearing. Its bodily
wounds therefore cannot be read as a disclosure that the narrator has
secretly judged him wicked. The chant compresses the two opening chapters:
permission concerning possessions, their destruction, and the later
attack on his flesh. The Missal's marginal ``Job~1'' is a broad locator;
the bodily attack belongs to chapter~2. The scriptural restrictions on
Satan's action remain essential to understanding the condensed narrative.

Gregory notices that Satan requests permission. The enemy's will is
malignant, but he possesses no independent sovereignty over creation.
Gregory also asks why God can say Job was afflicted ``without cause.''
His answer distinguishes an antecedent crime deserving punishment from
the manifestation and growth of virtue. Job's trial lacks the former;
it does not lack every providential end. These distinctions answer a
question raised by the narrative; they do not provide an observer with
God's reason for a particular neighbor's calamity.\footnote{Gregory the
Great, \work{Moralia in Iob} II.10.16--11.19; III.3.3--6.8, historical
Library of Fathers English (1844), retained transcription.}

\subsection{Restoration through the offering (Sunday Secret)}
The Secret asks God to receive the offerings favorably, and names both
propitiation---the restoration of peace with God through sacrifice---and
restored salvation as effects he himself has willed.
The phrase \textit{potenti pietate} unites power with loving goodness.
The community does not approach an unwilling God with a price it has
independently devised. It appeals to the divine saving will that gives
the offering its meaning. Schuster develops the double movement of
reconciliation and restoration through the Eucharistic sacrifice.\footnote{
Schuster, \work{Sacramentary} III, p.~181.}

Job's appearance immediately before this prayer gives suffering a place
within the assembly's offering, but Job and the sacrificed Christ are
not interchangeable figures. Schuster reads Job typologically as Christ
innocent and suffering for our sins. The innocent Christ's redemptive
work supplies the fulfillment; the literal story still rebukes the
friends' inference that Job's pain proves his own wickedness.

\subsection{One God, equal majesty (Trinity Preface)}
The Preface moves from perpetual thanksgiving to a precise confession:
one God and one Lord, with distinction of Persons, unity of essence,
and equality of majesty. Divine mercy, power, and judgment are not
separate divine agendas. In particular, the Father is not assigned a
justice from which a kinder but lesser Son rescues us. Augustine's
account of the divine missions makes the relevant distinction: the Son's
being sent does not make him less than the Father, nor the Spirit's
being sent less than either.\footnote{Augustine, \work{On the Trinity},
III, preface \S3, Arthur West Haddan historical English as delivered by
New Advent. This is doctrinal illumination of the Preface.}

The angels' praise leads into the Sanctus of the Ordinary. People engaged
in the Epistle's conflict already join that worship. Their thanksgiving
does not wait until all opposition has disappeared.

\subsection{Hope asks when (Communion)}
The chant selects Psalm~118:81,84,86. It leaves out the Bible's opening
verb of fainting, the question about the servant's days, and other
clauses; it retains hope in the word, the question about judgment, the
report of persecution, and the cry for help. Augustine's treatment of
the full psalm supplies context without replacing these appointed
words. His distinction is exact: the soul fails \emph{toward} salvation
in longing, not \emph{away} from it in despair.\footnote{Augustine,
Ps.~118, Caph, on vv.~81--88, especially vv.~81--82.}

The demand for judgment stands beside the Gospel's demand for pardon.
Augustine hears the martyrs' question in the Church's voice; Bellarmine
hears desire for lasting peace rather than curiosity about the number
of days left in the world. Schuster places another emphasis beside
them: rejection of truth already carries a judgment in the loss of
truth's light. That present judgment does not exhaust the psalmists'
future expectation. The question ``when?'' remains a real prayer.

\subsection{The mouth and the mind (Sunday Postcommunion)}
The Eucharist is named \textit{immortalitatis alimoniam}, nourishment
of immortality. The prayer then asks that the mind follow what the mouth
has received. Bodily reception matters, but it is not the endpoint of
the petition. Schuster develops its claim upon mind, heart, and will.
The inherited English says ``receive with pure souls''; the Latin
\textit{sectemur} supplies the further force of following or pursuing
in the commentary. The prayer's human translation remains as it
stands.\footnote{Cummiskey, XXI Sunday after Pentecost, Postcommunion;
Schuster, III, p.~181.}

\subsection{A gift sent among the nations (mission additions)}
When the mission orations are appointed, the Collect's request that
all be saved is expressed through workers sent into the harvest and
given confidence to speak. The prayer joins the saving will of
1 Timothy~2, the progress of the word in 2 Thessalonians~3, and the
knowledge of the Father and the sent Christ in John~17. These are
scriptural resonances within a composed prayer, not extra readings
added to the Sunday.

The mission Secret names Christ who gave himself as redemption for all,
and asks that God's name be magnified from sunrise to sunset through a
pure offering. Its geographical breadth challenges a household enclosed
in itself. The Postcommunion returns to the people already nourished:
true faith must continue to advance through redemption's gift. The
sequence is sending, universal worship, and enduring growth. The
historical Lasance English makes some freer choices, noted with the
texts; the Latin's demand that faith remain true is indispensable to
the argument.

\subsection{Service marked by the cross (St Luke commemoration)}
Where Luke is commemorated, the Collect speaks of bearing the cross's
mortification in his body for God's name. In 2 Corinthians~4:10--18,
bodily affliction and the manifestation of Jesus' life belong together;
ministry directs attention to Christ, not to the minister's self-display.
The prayer applies that apostolic pattern to the Evangelist. It does
not by itself establish the historical manner of his death.

The Secret joins free-minded service to healing and glory, and the
Postcommunion seeks the sanctification and safety of the soul through
the altar's gift and the saint's prayer. These requests add a personal
witness to the Sunday's household. Freedom serves God; healing opens
toward glory; safety is prayed for by people whose saint is remembered
as a bearer of the cross. No separate Luke Gospel or Apostles Preface
enters this Sunday's appointment.
